\subsection*{Objectifs du TP}

Vous allez construire un \textbf{mini-jeu de combat au tour par tour} (façon RPG) en langage C,
pour mettre en pratique les types structurés : structures simples, structures imbriquées,
tableaux de structures, et énumérations.

\textbf{Compétences visées :}
\begin{itemize}
  \item Déclarer et utiliser des structures simples et imbriquées
  \item Manipuler des tableaux de structures
  \item Passer des structures en paramètres de fonctions, et les modifier en les renvoyant
    avec \texttt{return}
  \item Utiliser un type énuméré (\texttt{enum}) comme champ d'une structure
  \item Combiner ces notions pour résoudre un problème plus riche (mini-projet final)
\end{itemize}

\subsection*{Organisation}

Vous travaillez sur un \textbf{unique fichier}, \texttt{jeu\_combat.c}, du début à la fin de la
séance. Vous ne créez pas de nouveau fichier à chaque exercice : vous partez du squelette
fourni, et chaque étape du TP s'ajoute au code déjà écrit. Vous ne connaissez pas encore les
pointeurs (\texttt{*}, \texttt{\&}, \texttt{malloc}) : ce TP n'en a pas besoin.

\textbf{Durée estimée :} \UPSTIduree

\begin{UPSTIinfor}{Consigne importante}
  Le squelette \texttt{jeu\_combat.c} est organisé en \textbf{quatre sections fixes}, repérées
  par des commentaires : déclaration des structures, déclaration des fonctions (prototypes),
  programme principal (\texttt{main}), définition des fonctions (corps). À chaque étape,
  rangez bien chaque nouveau morceau de code dans la section correspondante : une structure
  déclarée avec les structures, un prototype avec les prototypes, le code du jeu dans
  \texttt{main}, le corps d'une fonction avec les corps de fonctions. Ce rangement n'est pas
  une question de style : il garantit que votre code compile (l'ordre de déclaration compte en
  C) et qu'il reste lisible du début à la fin de la séance.
\end{UPSTIinfor}

\pagebreak
